\documentclass[12pt]{article}
\usepackage[margin=1in]{geometry}
\usepackage[T1]{fontenc}
\usepackage{bold-extra}
\usepackage{xcolor}
\definecolor{garnet}{HTML}{73000A}

\setlength{\parindent}{0pt}
\setlength{\parskip}{0.5em}

\begin{document}

{\color{garnet}\large\textbf{\textsc{SC Space Grant Personal Essay}}}\\[-0.8em]
{\color{garnet}\rule{\textwidth}{0.5pt}}\\[0.2em]
The first time I fully appreciated the complexity of how aerospace engineers interact with the physical world was during my Honors Introduction to Aerospace Engineering course. As the class integrated concepts from statics, controls, system dynamics, and aerodynamics, I saw how aerospace engineering requires balancing theoretical models with practical data and constraints. This experience shifted my interest from a general fascination with aerospace systems to a focused curiosity about developing effective representations of complex systems that enable meaningful analysis and innovation.

Following this shift, I developed a simplified simulation of aircraft powertrains using MATLAB and Simscape to better understand how complex systems can be represented efficiently without sacrificing physical relevance. In parallel, I explored aircraft modeling using computer-aided design software, which further strengthened my understanding of system structure and component interaction. These projects emphasized the importance of selecting appropriate assumptions and demonstrated that effective simplification requires a clear understanding of practical constraints.

In parallel, my involvement with the USC Fly Club provided valuable hands-on context. Discussions with pilots offered firsthand insight into how aircraft systems function in practice and how interconnected components can often be replaced or approximated by simpler mechanisms. One particularly illustrative example came from a gyrocopter pilot who constructed his own aircraft using components adapted from cars, helicopters, airplanes, and snowmobiles.

These experiences directly informed my simulation work and motivated me to apply similar modeling principles to subsequent projects, including the ongoing design of a hovercraft and adjustments to traditional configurations aimed at improving efficiency.

Together, these experiences have solidified my interest in developing simplified, physics-based models that connect theoretical analysis with real-world system behavior. As I advance in my studies, I am increasingly applying these approaches to mechanical and aerospace systems to evaluate design trade-offs and guide early-stage improvement. Through continued research experiences, I aim to build the technical foundation necessary to contribute to advanced aerospace systems and space exploration efforts.

Recent advances by Boom Supersonic and NASA's X-59 QUESST program illustrate how progress in aerospace increasingly depends on reexamining long-standing assumptions and improving how complex systems are modeled and validated. Preparing to contribute to this future motivates my focus on developing adaptable, physically grounded models that support innovation under realistic constraints.

Support from the NASA South Carolina Space Grant would allow me to further refine these modeling and systems-level skills through research-driven projects. By providing the time, mentorship, and resources necessary for focused exploration, Space Grant would prepare me to contribute meaningfully to advanced aerospace and space exploration efforts, where effective system representation and validation are essential to innovation.

\end{document}
